\chapter{Introduction}

\section{Overview}

A compiler translates a program written in a high-level language into a lower-level form that a machine can execute. Between the front end, which reads and checks the source program, and the back end, which emits machine instructions, sits the optimizer. The optimizer rewrites an intermediate representation (IR) of the program so that the final code is smaller, faster, or both, while preserving the meaning of the original program exactly. Classical textbooks describe the optimizer as a sequence of \emph{passes}, each of which performs one well-defined transformation: folding constant expressions, deleting code whose results are never used, reusing values that have already been computed, replacing expensive operations by cheaper ones, and so on \cite{aho2006,cooper2011,muchnick1997}.

Each of these passes is easy to understand in isolation. What is not easy is understanding how they behave together. Passes are not independent: the output of one pass is the input to the next, and a transformation can create new opportunities for a later pass, or remove opportunities that a later pass would otherwise have exploited, or leave behind residue that only a specific later pass can clean up. Constant folding, for example, turns \code{t1 = 2 + 3} into \code{t1 = 5}. That, by itself, is a small improvement. Its real value is indirect: once \code{t1} is a known constant, every instruction that reads \code{t1} can be simplified, comparisons that depend on \code{t1} may become decidable at compile time, whole branches may become unreachable, and the code that computed the operands of the original addition becomes dead. None of those follow-on improvements happen unless the passes that perform them run \emph{after} constant folding. If dead code elimination runs first, it finds very little to remove, and it does not run again.

The question of which order to run the passes in is known as the \emph{phase-ordering problem}. It is one of the oldest open problems in compiler construction. It is hard for three reasons. First, the number of possible orderings grows factorially: six passes have $6! = 720$ orderings, ten passes have more than three million, and a production compiler with a hundred or more passes has a space that can never be enumerated. Second, the best ordering is program dependent: a sequence that is excellent for a loop-heavy numerical kernel may be mediocre for a branchy state machine. Third, the interactions between passes are difficult to reason about analytically, because each pass is a complex procedure over an arbitrary program, and formal results show that the problem is undecidable in its general form \cite{touati2006}.

Production compilers such as GCC and LLVM \cite{lattner2004} resolve the problem pragmatically. Their optimization levels, \code{-O1}, \code{-O2}, \code{-O3}, and \code{-Os}, are fixed pipelines in which the pass order was settled by decades of engineering experience and tuning against large benchmark suites. This approach works well on average, and it is the right approach for a compiler that must serve millions of users. It has two costs, however. The first is that the pipelines are opaque: a student, a researcher, or an embedded-systems engineer who wants to know \emph{why} a given order was chosen, or how much a different order would gain or lose on their particular code, has no direct way to find out. The second is that the pipelines are tuned for typical desktop and server workloads; the trade-offs that matter on a small microcontroller, where every byte of flash and every cycle of energy counts, may be different.

This project approaches the phase-ordering problem from the opposite direction. Instead of studying a production compiler from the outside, the team built a complete compiler infrastructure from scratch in which every optimization pass is an independent, reorderable unit, and in which the space of orderings is small enough to enumerate completely. The compiler accepts a subset of the C language and is implemented in Python. It comprises a hand-written lexer, a recursive-descent parser, a semantic analyzer, a three-address-code intermediate representation, six optimization passes managed by a configurable pass manager, an interpreter that executes the intermediate code to measure dynamic behaviour and to check correctness, and a back end that emits RISC-V RV32IM assembly. Fifteen benchmark kernels adapted from the MiBench embedded suite \cite{guthaus2001} and the PolyBench/C numerical suite \cite{pouchet2016} are compiled under all 721 possible orderings of the six passes, and three metrics are collected for each result: static code size, an estimated cycle count from an embedded-processor cost model, and the dynamic instruction count observed during execution.

The infrastructure is complemented by an interactive web application. Its purpose is to make the phase-ordering problem \emph{visible}: a user can type or load a program, see the tokens, the syntax tree, the symbol table, and the intermediate code, arrange the six passes in any order by drag and drop, step through the pipeline one pass at a time with a coloured diff of each stage, inspect the generated RISC-V assembly side by side with the intermediate code, and run the exhaustive 721-ordering exploration directly in the browser with an interactive plot of the size-versus-cycles trade-off. The tool is intended for compiler courses and for anyone who wants to build intuition about pass interactions without reading the source of a production compiler.

\begin{figure}[H]
\centering
\resizebox{\textwidth}{!}{%
\begin{tikzpicture}[
  node distance=4mm and 6mm,
  box/.style={draw, rounded corners=2pt, minimum width=1.55cm, minimum height=0.8cm, align=center, font=\footnotesize\sffamily, fill=white},
  ir/.style={draw, rounded corners=2pt, minimum width=2.1cm, minimum height=0.8cm, align=center, font=\footnotesize, fill=gray!8},
  res/.style={draw, rounded corners=2pt, minimum width=2.1cm, minimum height=0.8cm, align=center, font=\footnotesize, fill=blue!8},
  arr/.style={-{Stealth[length=2.2mm]}, thick},
  lab/.style={font=\scriptsize\itshape, text=black!60}]
  \node[ir] (in1) {factorial IR\\40 instructions};
  \node[box, right=of in1] (a1) {CF};
  \node[box, right=of a1] (a2) {CSE};
  \node[box, right=of a2] (a3) {CP};
  \node[box, right=of a3] (a4) {DCE};
  \node[box, right=of a4] (a5) {SR};
  \node[box, right=of a5] (a6) {AS};
  \node[res, right=of a6] (o1) {16 instructions\\$-60.0\%$};
  \foreach \s/\t in {in1/a1,a1/a2,a2/a3,a3/a4,a4/a5,a5/a6,a6/o1} \draw[arr] (\s) -- (\t);

  \node[ir, below=14mm of in1] (in2) {factorial IR\\40 instructions};
  \node[box, right=of in2] (b1) {DCE};
  \node[box, right=of b1] (b2) {CF};
  \node[box, right=of b2] (b3) {CSE};
  \node[box, right=of b3] (b4) {CP};
  \node[box, right=of b4] (b5) {SR};
  \node[box, right=of b5] (b6) {AS};
  \node[res, right=of b6, fill=orange!12] (o2) {33 instructions\\$-17.5\%$};
  \foreach \s/\t in {in2/b1,b1/b2,b2/b3,b3/b4,b4/b5,b5/b6,b6/o2} \draw[arr] (\s) -- (\t);

  \node[lab, above=1mm of a1.north west, anchor=south west] {Ordering A: enabling passes first, clean-up last};
  \node[lab, above=1mm of b1.north west, anchor=south west] {Ordering B: same six passes, dead code elimination first};
\end{tikzpicture}}
\caption[The phase-ordering problem in one picture.]{The phase-ordering problem in one picture. The same six passes applied to the same intermediate code of the \code{factorial} benchmark produce a 16-instruction program in one order and a 33-instruction program in another. The only difference is where dead code elimination runs.}
\label{fig:intro-phase-ordering}
\end{figure}

\fig{fig:intro-phase-ordering} shows the effect that motivates the entire project. The unoptimized intermediate code of the \code{factorial} benchmark has 40 instructions. Running the six passes in the order constant folding, common subexpression elimination, copy propagation, dead code elimination, strength reduction, algebraic simplification yields 16 instructions, a reduction of 60\%. Moving dead code elimination to the front and leaving every other pass in the same relative order yields 33 instructions, a reduction of only 17.5\%. No pass was added or removed. The interaction between the passes, and nothing else, accounts for a factor of two in the outcome.

\section{Scope}

The project has a deliberately bounded scope so that the phase-ordering question can be answered completely rather than approximately. The following are within scope and were delivered.

\begin{itemize}
  \item \textbf{Source language.} A subset of C consisting of the \code{int} type, integer literals, arithmetic operators (\code{+}, \code{-}, \code{*}, \code{/}, \code{\%}), comparison operators, logical operators with C short-circuit semantics, unary negation and logical not, global scalar variables with constant initializers, fixed-size one-dimensional arrays, block scoping with shadowing, \code{if}/\code{else}/\code{else if}, \code{while} and \code{for} loops, functions with parameters and return values, a built-in \code{print} statement, and single-line and block comments. Everything is a 32-bit integer; there is no floating point, no pointers, no structures, and no standard library.
  \item \textbf{Compiler front end.} A lexer with line and column tracking, a recursive-descent parser with six levels of operator precedence, an abstract syntax tree of twenty node types, a scope-aware symbol table, and a semantic analyzer that collects every violation in a program and reports them together.
  \item \textbf{Intermediate representation.} A flat three-address code with 35 opcodes, including dedicated opcodes for array and global-variable access so that block-local optimizations cannot miscompile cross-function state.
  \item \textbf{Optimization passes.} Six classical passes: constant folding with constant propagation (\pass{CF}), dead code elimination with unreachable-code removal (\pass{DCE}), common subexpression elimination (\pass{CSE}), copy propagation (\pass{CP}), strength reduction (\pass{SR}), and algebraic simplification (\pass{AS}). Each pass is a pure function from an instruction list to a new instruction list, and each is safe to run in any position of any ordering.
  \item \textbf{Pass manager.} A component that validates and runs an arbitrary ordering, describes it, and enumerates all 720 full permutations plus the baseline.
  \item \textbf{Measurement.} A metric collector that computes static code size, a weighted cycle estimate based on an ARM Cortex-M class cost model, and an instruction-category breakdown; an interpreter that executes the intermediate code and reports dynamic instruction counts, per-opcode frequencies, and program output; and a visualizer that produces seven kinds of publication-quality plots and a geometric-mean summary.
  \item \textbf{Back end.} A RISC-V RV32IM assembly generator producing GNU-syntax, Linux user-mode assembly with a built-in decimal print routine, validated against an in-repository RV32IM simulator so that the assembly is checked to match the interpreter on every benchmark.
  \item \textbf{Benchmarks.} Fifteen kernels: eight adapted from MiBench (bit counting, Collatz iteration, factorial, Fibonacci, greatest common divisor, integer square root, modular exponentiation, and a SHA-style mixing function) and seven adapted from PolyBench/C (one- and two-dimensional Jacobi stencils, GEMM, 2MM, ATAX, BiCG, and GESUMMV).
  \item \textbf{Interactive tool.} A React and TypeScript single-page application served by a FastAPI back end, with five sections: Playground, Optimization Lab, Assembly, Phase-Ordering Explorer, and Reference. A zero-dependency legacy page is also served for environments without a JavaScript build.
  \item \textbf{Quality assurance.} A test suite of 185 automated tests covering every module, including a cross-pass integration test that compiles all fifteen benchmarks under every ordering and asserts output equality with the baseline.
\end{itemize}

The following are explicitly outside the scope of this phase. The compiler does not perform register allocation, instruction scheduling, loop transformations such as unrolling or invariant code motion, inlining, or interprocedural analysis. The optimization passes use local, forward, linear analyses that conservatively reset their state at control-flow boundaries rather than full data-flow analysis over a control-flow graph. The cycle estimate is a static cost model, not a measurement on hardware or in a cycle-accurate simulator. The RISC-V back end uses a simple stack-slot allocation strategy and is intended for faithful, readable lowering rather than peak performance. These exclusions are not weaknesses of the study; they are what make the study tractable. The goal was to characterize the interactions of a small, well-understood set of passes completely, and to build the infrastructure on which more ambitious passes can later be added.

\section{Motivation}

The motivation for this work comes from three directions: education, research, and the practical needs of embedded software.

\textbf{Education.} Compiler design is a core subject in every computer science curriculum, and optimization is usually taught pass by pass: one lecture on constant folding, one on dead code elimination, one on common subexpression elimination. Students learn what each pass does, but they rarely see what the passes do to \emph{each other}. The reason is practical: production compilers are enormous, their pass pipelines are buried under layers of infrastructure, and their intermediate representations are hard to read. Online tools such as Compiler Explorer \cite{godbolt2019} show the final assembly for a given optimization level but do not expose the pipeline that produced it. Recent educational work has begun to address this gap with visualizations of individual optimizations \cite{koitz2025}, but a tool in which a learner can freely rearrange passes and watch the consequences unfold stage by stage did not exist. The team wanted to build one, and to build it on a compiler small enough that a student can read all of its source code in an afternoon.

\textbf{Research.} The phase-ordering problem has a rich literature. Early work used genetic algorithms and other search methods to find good sequences \cite{cooper1999,almagor2004,kulkarni2004}. Kulkarni and colleagues showed that for small sets of passes the space of \emph{distinct outcomes} is far smaller than the space of sequences, because many sequences produce identical code, and that this makes exhaustive enumeration feasible for realistic functions \cite{kulkarni2006,kulkarni2009}. More recent work applies machine learning and reinforcement learning to predict good orderings for LLVM \cite{ashouri2017,hajali2020,jain2022,cummins2023}. Almost all of this work is conducted on production compilers, which means the researcher can never be sure that a measured effect is due to pass interaction rather than to an artefact of the surrounding infrastructure. A small, transparent compiler with a handful of passes offers a complementary view: every ordering can be run, every result can be inspected by hand, and the structural rules that govern the space can be stated exactly. The project is an attempt to produce such a view for six classical passes on embedded-style kernels.

\textbf{Embedded software.} On desktop processors, code size is rarely a concern and the compiler's default pipeline is nearly always good enough. On microcontrollers it is different. A device built around a Cortex-M0 class core may have 32 kilobytes of flash; the difference between a 60\% and a 20\% reduction in code size can decide whether a feature fits at all. Energy consumption is tied directly to the number of cycles executed, so a 20\% reduction in dynamic instructions translates into longer battery life. Studies on embedded targets have found that the standard optimization levels are not always the best choice and that alternative sequences can improve both performance and energy \cite{georgiou2018}. Small manufacturers who build such devices, including the large population of micro, small, and medium enterprises in India, usually cannot afford to tune a compiler for their workload. A framework that quantifies how much ordering matters, and that identifies the structural rules for a good ordering, is a first step towards a tool that could do such tuning automatically. The choice of MiBench, which was designed to represent commercial embedded workloads, and of a RISC-V RV32IM back end, which targets an open instruction set that is increasingly used in low-cost microcontrollers, reflects this motivation.

\textbf{Continuity from Phase I.} The first phase of this project delivered the front end, the intermediate representation, three passes (constant folding, dead code elimination, and common subexpression elimination), a first pass manager, eight MiBench-adapted benchmarks, and a single-file web page for experimentation. The results of that phase, obtained on the \code{factorial} benchmark, showed a 60\% code-size reduction for the best ordering and confirmed that the effect was real and large. Phase II set out to answer the questions that Phase I raised: does the effect hold across a broader and more realistic set of programs, in particular array-based numerical kernels; are there structural rules that explain which orderings are good; how much of the ordering space is redundant; does the improvement in the intermediate code survive lowering to a real instruction set; and can the whole apparatus be made usable by others? Answering those questions required three more passes, seven more benchmarks, a semantic analyzer, arrays, global variables, \code{for} loops, short-circuit evaluation, a RISC-V back end with a simulator to validate it, a redesigned web application, and a substantially larger test suite. This report describes that work and its results.

\textbf{Organization of the report.} Chapter 2 surveys the literature on the phase-ordering problem, on compiler autotuning, on the benchmark suites and evaluation methodology used, and on educational compiler tools. Chapter 3 formulates the problem, states the objectives, and lists the functional and non-functional requirements and the technical specifications. Chapter 4 presents the architecture, high-level and detailed design of every component, describes the fifteen benchmark programs, and lists the tools and technologies used. Chapter 5 reports the experimental environment, the results achieved against each requirement, the empirical findings of the exhaustive phase-ordering study, a comparison with existing work, snapshots of the interactive tool, and the societal impact and Sustainable Development Goals addressed. Chapter 6 concludes and outlines future work. The appendices contain the grammar of the source language, the intermediate-representation opcode reference, a complete compilation trace for one benchmark, and a summary of the test suite.
